\seqsplit{ClassificacaoDoRegistroDeReceitaOrcamentaria} &
Relator que associa a Receita às classificações de natureza, fonte e conta contábil. 
\\ \hline

\seqsplit{ContaContabil} &
Conta é a expressão qualitativa e quantitativa de fatos de mesma natureza, que evidencia a composição, variação e estado do patrimônio, bem como de bens, direitos, obrigações e situações nele não compreendidas, mas que, direta ou indiretamente, possam vir a afetá-lo. Cite MCASP - 11º Edicao: Parte IV: Seções 2.3 ; 
\\ \hline

\seqsplit{FonteDaReceita} &
No momento da arrecadação, a FR deverá ser informada no registro da realização da receita orçamentária (classe 6 do PCASP), no registro do ingresso dos recursos financeiros no ativo financeiro (classe 1 com atributo "F") e em contas de controle da Disponibilidade por Destinação de Recursos - DDR (nas classes 7 e 8 do PCASP). A partir desse momento, a FR informada para os recursos financeiros não deverá ser alterada Cite MCASP - 11º Edicao: Parte I: Seção 5.2 
\\ \hline

\seqsplit{NaturezaDaReceita} &
O § 1º do art. 8º da Lei nº 4.320/1964 define que os itens da discriminação da receita, mencionados no art. 11 dessa lei, serão identificados por números de código decimal. Convencionou-se denominar este código de natureza de receita. Importante destacar que a classificação da receita orçamentária por natureza é utilizada por todos os entes da Federação e visa identificar a origem do recurso segundo o fato gerador: acontecimento real que ocasionou o ingresso da receita nos cofres públicos. A classificação por Natureza de Receita Orçamentária é composta por um código de oito dígitos numéricos que representam: a Categoria Econômica, a Origem, a Espécie, os Desdobramentos e o Tipo de Receita. Cite MCASP - 11º Edicao: Parte I: Seção 3.2.1 
\\ \hline

\seqsplit{ReceitaOrcamentaria} &
O orçamento é um importante instrumento de planejamento de qualquer entidade, seja pública ou privada, e representa o fluxo previsto de ingressos e de aplicações de recursos em determinado período. Cite MCASP - 11º Edicao: Parte I: Seção 3.1 São disponibilidades de recursos financeiros que ingressam durante o exercício e que aumentam o saldo financeiro da instituição. Instrumento por meio do qual se viabiliza a execução das políticas públicas, as receitas orçamentárias são fontes de recursos utilizadas pelo Estado em programas e ações cuja finalidade precípua é atender às necessidades públicas e demandas da sociedade. Essas receitas pertencem ao Estado, transitam pelo patrimônio do Poder Público e, via de regra, por força do princípio orçamentário da universalidade, estão previstas na Lei Orçamentária Anual - LOA. Cite MCASP - 11º Edicao: Parte I: Seção 3.1.2 
\\ \hline

\seqsplit{ReceitaOrcamentariaRegistrada} &
Receita orçamentária registrada emitida por uma unidade gestora e beneficiando outra unidade gestora. 
\\ \hline

\seqsplit{RegistroDeReceitaOrcamentaria} &
Registro que associa a Receita às unidades gestoras emitentes e beneficiadas. 
\\ \hline

\seqsplit{UnidadeGestoraBeneficiadaDeReceita} &
Unidade gestora beneficiada por uma receita orçamentária emitida. 
\\ \hline

\seqsplit{UnidadeGestoraEmitenteDeReceita} &
Unidade gestora responsável pela emissão da receita orçamentária. 
\\ \hline

\seqsplit{ValorAtreladoNaReceita} &
Tipo de dados que agrupa os valores monetários que um Registro de Receita Orçamentária pode ter. 
\\ \hline
